\section{Introduction}
\label{sec:introduction}
\subsection{Relevance}
\label{sec:intro:relevance}

Energy system modeling undergoes fundamental transformations, driven by
global developments in the field of Artificial Intelligence (AI) and the
growing need to model sector-coupled multi-energy systems (MES) at high
spatio-temporal resolution \cite{Klemm2021,Lim2025}.
Extensive research on AI-related methods has
been conducted to address the computational challenges of decentralized
intermittent energy production, the coupling of the electricity, heat,
mobility, and industry sectors with emerging technologies such as
hydrogen, power-to-X, and electric vehicles (EVs), and multi-objective
optimization (MOO), where economic, environmental, technical, and
resilience-related objectives are considered simultaneously
\cite{Zhou20181,Zheng20231907,Ren2024,Dong2023,Yin2024}. These
developments highlight the need for novel modeling approaches to capture
the multi-dimensionality of future energy systems and its implications
towards high-dimensional, stochastic, and nonlinear optimization models
while remaining computationally feasible for efficient and sustainable
energy policy decision-support \cite{Ruan2021221,Perera2019191,Khaloie2025,Mylonopoulos202332697}.

Surrogate models - also referred to by different authors as metamodels
\cite{Roth2019214,Hirvonen2017323,Jia2022553,Xiao2018,Li2023,Wang2020202933,Starke2025214},
response surfaces
\cite{Ahmad202451,Salahinezhad2025,Xu20212696,Schulte2020,Roy2023},
emulators \cite{Hu2020,Volodina2022,Lawson201647}, kriging models,
Gaussian-process regressions
\cite{Deng2020831,Nikolaidis2021,Li2021,Perera2019191}, or
polynomial-chaos expansion models
\cite{Muhlpfordt20194806,Safta20172535,Koirala20242284,Xu20212696} -
can be defined as a data-driven or statistical approach that replaces
repeated evaluations of a computationally expensive optimization model
with a faster approximation
\cite{Perera2019191,Ruan2021221,Starke2025214,Khaloie2025,Mylonopoulos202332697,Su2022315,Thrampoulidis2021,Han2018622}.
A set of representative input--output samples is first generated by
running the original model for selected system configurations. Based on
these samples, a machine-learning model is trained to approximate
decision variables, system parameters, and relevant model outputs.

Reviews of learning-assisted power system optimization have shown that
mixed-integer dispatch, AC optimal power flow, and capacity expansion
planning routinely become intractable when stochastic recourse, physical
network detail and rolling-horizon control are added
\cite{Ruan2021221,Khaloie2025,Lim2025,Mylonopoulos202332697,Perera2019191}.
These reviews also highlight surrogate modeling to be particularly
powerful in achieving substantial reduction in computational costs while
maintaining sufficient predictive accuracy
\cite{Ruan2021221,Perera2019191,Starke2025214,Khaloie2025,Mylonopoulos202332697}.
Machine-learning methods have become routinely available, with
established libraries for kriging, polynomial chaos, random forests,
gradient boosting and deep neural networks
\cite{Perera2019191,Starke2025214,Khaloie2025,Elsheikh2019622}.
Several recent reviews have addressed sub-questions in this regard:
Perera et al.Â~\cite{Perera2019191} reviewed machine-learning methods in
energy system optimization with a strong emphasis on building and
microgrid applications. Ruan et al.Â~\cite{Ruan2021221} reviewed
learning-assisted power system optimization and Khaloie et
al.Â~\cite{Khaloie2025} focused on machine-learning techniques for optimal
power flow. Mylonopoulos et al.Â~\cite{Mylonopoulos202332697} covered
optimization methods specifically in ship energy systems, and Lim et
al.Â~\cite{Lim2025} reviewed power system decision making in the
deep-learning era. Starke et al.Â~\cite{Starke2025214} reviewed the
applicability of machine learning to the metamodeling of energy systems.
The present review consolidates and extends these contributions in four
ways.

\subsection{Scope and contribution of this review}
\label{sec:intro:scope}

We focus on optimization-based surrogate modeling in energy system
modeling, specifically on MES modeling and MOO. The review covers the
main optimization tasks in which surrogates have been used in the last
decade, including economic dispatch and unit commitment
\cite{Wu2022231,Chen20244723,Chen2022,Nikolaidis2021,Lu20232989,Safta20172535,Wang20211767,Pang2023},
optimal power flow and chance-constrained variants
\cite{Muhlpfordt20194806,Su2022315,Deng2020831,Liu20223726,Chen2023657,Koirala20242284,Xu20212696,Srithapon202134395,Urgun20202791},
capacity and generation expansion
\cite{Rodgers2018951,Jahangiri20244849,Jahangiri2023,Han2018622,Fu20202904,Fu2022},
district heating and integrated heating-cooling systems
\cite{Koschwitz2018134,Zhou2024,Jia2022553,Hirvonen2017323,Reynolds2017704,Gao2025,Petrova2025},
multi-energy and sector-coupled systems
\cite{Zhou20181,Zheng20231907,Ren2024,Jalving2023,Dong2023,Yin2024},
microgrid and energy-hub planning and operation
\cite{Faraji2020157284,Nematollahi2021,Wang20192635,Hussien20211883,Rigo-Mariani20171762,Cao2023__2,Li2023__3},
multi-objective design
\cite{Thrampoulidis2021,Hussien20211883,Schulte2020,Han2018622},
and stochastic planning
\cite{Zhou20181,Fu20202904,Fu2022,Yang2024,Dong2022,Chen2023657,Xu20212696,Liang20246508,Dong2023}.

This review makes four contributions. First, we integrate the literature
at the method level by providing a cross-cutting taxonomy of the
surrogate models used in energy system optimization, together with their
optimizer-compatibility properties, such as gradient-friendliness,
convexity and mixed-integer compatibility, and their typical failure
modes. Second, we identify five recurring integration patterns and link
them to the dominant energy-system tasks: P1 the surrogate replaces an
expensive subroutine, P2 the surrogate accelerates the search, P3 the
surrogate acts as a warm-start or solution proxy, P4 the surrogate
enables decomposition, and P5 the surrogate handles uncertainty. This
five-pattern view shows that the same surrogate model can play very
different roles across model tasks and explains why similar model
architectures appear in studies that are otherwise only loosely
connected. Third, we argue that the dominant validation practice of
reporting RMSE, MAE or R$^2$ is insufficient when the surrogate is
embedded in an optimizer, and we summarize the decision-aware metrics
that should complement predictive accuracy metrics. Fourth, we publish a
machine-readable evidence map and the accompanying
\texttt{surrogates_esm_screening_enriched.csv} that links each reviewed study
to a task class, surrogate family, integration pattern and reported
decision impact. The evidence map is intended to support filtering,
comparison and reuse in follow-up work.

The remainder of the paper is structured as follows:
SectionÂ~\ref{sec:optimization-bottlenecks} establishes the optimization-side
vocabulary, introduces the dominant problem classes in energy system
modelling and explains where the computational load arises.
SectionÂ~\ref{sec:taxonomy} presents the surrogate-side taxonomy.
SectionÂ~\ref{sec:doe}
discusses training data and design of experiments.
SectionÂ~\ref{sec:integration} develops the five integration patterns and
maps them onto problem classes.
SectionÂ~\ref{sec:validation} addresses validation and decision-aware
metrics. SectionÂ~\ref{sec:applications} consolidates the literature by application
domain and synthesizes when surrogates demonstrably pay off.
SectionÂ~\ref{sec:challenges} formulates open challenges and a research
agenda. SectionÂ~\ref{sec:conclusion} concludes.
